\section{はじめに}
画像データセットに対する物体アノテーションは，画像中のどこに物体が存在するかという物体領域の判断と，その物体が何であるかというカテゴリの判断の2種類から構成される．
You Only Look Once（YOLO）\cite{ref_YOLO}やDETR\cite{ref_DETR}に代表される一般的な物体検出モデルは，これら2つの判断を単一のモデルに統合し，あらかじめ定めた固定のクラス集合に対して同時に学習する．
そのため，学習時に定めたクラス集合の範囲内では高精度な検出が可能である一方，学習後に新たなカテゴリを追加したり，カテゴリの粒度を変更したりする場合には，モデル自体の再学習が必要になる．

近年では，OWL-ViT\cite{ref_OWLViT}やGrounding DINO\cite{ref_GDINO}のように，カテゴリ名を自然言語のテキストクエリとして与えることで，固定のクラス集合に依存せずに物体を検出するオープンボキャブラリ物体検出モデルが提案されている．
しかし，これらのモデルも物体領域の検出とカテゴリの予測を単一のモデルの中で同時に行う．
そのため，例えば「動物」という粗いカテゴリで検出した物体を「犬」「猫」のようなより細かいカテゴリで改めて分類したい場合，候補カテゴリ集合を変更した上で検出処理そのものを再度実行する必要がある．

本研究では，物体領域の検出とカテゴリのアノテーションという2つの判断を明確に分離したアーキテクチャを提案する．
物体領域の検出には，カテゴリに依存しないバウンディングボックスを1回の推論で出力するDETRを用いる．
カテゴリのアノテーションには，検出済みの領域を切り出した画像（クロップ）に対し，VLMによるゼロショット分類を適用する．
両者が分離されているため，同一の検出結果に対して候補となる語彙集合を変更するだけで，検出処理を再実行することなくカテゴリのアノテーションを繰り返すことができる．
本研究はこの性質を利用し，WordNet\cite{ref_WordNet}の上位語・下位語関係をたどることで，追加学習を行わずに「animal → dog → doberman」のような粒度可変のアノテーションを実現する手法を提案する．

本稿の貢献は以下の4点である．
第一に，物体領域の検出とカテゴリのアノテーションを分離するアーキテクチャと，その分離を利用したWordNetによる粒度可変アノテーションの手法を提案する．
第二に，COCO 2017 Validationを用いた定量評価により，カテゴリのアノテーションにはオープンボキャブラリなVLMが必要であること，および物体領域の検出にDETRを用いることが位置検出精度の観点から妥当であることを示す．
さらに，閉域物体検出モデルであるYOLOと同一の語彙で比較し，提案手法の利点と限界を定量的に示す．
また，YOLOの予測を深さ0に用い，深さ1以降の細分化のみをVLMで行うハイブリッド構成を評価し，YOLOの精度と本手法の粒度可変性を両立できることを示す．
\add{さらに，閾値の選択に用いない別の画像と，WordNetのSynsetが付与されたLVISのアノテーションを用いて，深さ0の誤りの伝播と深さ1以降の推定の正誤を評価し，累積確信度による細分化の停止規則を提案する．}
第三に，提案パイプライン全体のカテゴリアノテーション精度を評価し，論文等で明示されることの少ないクロップの前処理と検出スコアの閾値が精度に大きく影響することを示す．
第四に，WordNet階層の各深さにおけるVLMの確信度を定量的に評価し，粒度を細かくするにつれて確信度が変化する傾向と，WordNetの生の階層構造を候補語彙として用いることに起因する限界を，下位語が1語しかない分岐の頻度を含めて明らかにする．

本稿の構成は以下のとおりである．
2章では関連研究について述べる．
3章では提案手法について述べる．
4章では実験条件と結果を示し，5章で考察を行う．
6章で本稿をまとめる．
